\chapter{Dependencia continua respecto condiciones iniciales y parámetros}
\section{Flujo de sistemas de EDOS}
\begin{definition}[Flujo]
Sean $\displaystyle f : D \subset \R \times \R^{n} \to \R^{n} $ continua y localmente Lipschitz respecto de la segunda variable. Dado $\displaystyle \left(t_{0}, x_{0}\right)\in D $ definimos el \textbf{flujo} $\displaystyle \varphi\left(t; t_{0}, x_{0}\right) $ como la única solución del PVI con la condición inicial $\displaystyle \left(t_{0}, x_{0}\right) $. 
\end{definition}
\begin{observation}
	Tenemos que $\displaystyle \varphi : \tilde{D} \subset \R \times \R \times \R^{n} \to \R^{n} $ con 
	\[\tilde{D} = \left\{ \left(t; t_{0}, x_{0}\right) \; : \; \left(t_{0}, x_{0}\right)\in D, t \in \left(\beta_{-}\left(t_{0}, x_{0}\right), \beta_{+}\left(t_{0}, x_{0}\right)\right)\right\}  \]
	abierto.
\end{observation}
\begin{observation}
Por definición el flujo satisface
\[
\frac{\partial \varphi}{\partial t}\left(t; t_{0}, x_{0}\right)=f\left(t, \varphi\left(t; t_{0}, x_{0}\right)\right) \quad \text{y} \quad \varphi\left(t_{0}; t_{0}, x_{0}\right)=x_{0}
.\]
\end{observation}
\begin{lema}[Propiedad de semigrupo]
Para todo $\displaystyle \left(t_{0}, x_{0}\right)\in D $ y $\displaystyle s,t\in \left(\beta_{-}\left(t_{0}, x_{0}\right), \beta_{+ }\left(t_{0}, x_{0}\right)\right) $ se cumple que $\displaystyle \varphi\left(t; t_{0}, x_{0}\right)=\varphi\left(t; s, \varphi\left(s; t_{0}, x_{0}\right)\right) $. 
\end{lema}
\begin{proof}
Fijamos $\displaystyle \left(t_{0}, x_{0}\right)\in D $ y $\displaystyle s \in \left(\beta_{- }, \beta_{+ }\right) $. Definimos
\[y\left(t\right):=\varphi\left(t; t_{0}, x_{0}\right)\quad \text{y} \quad z\left(t\right):=\varphi\left(t; s, \varphi\left(s; t_{0}, x_{0}\right)\right) .\]
Tomamos también $\displaystyle x^{*} : = \varphi\left(s; t_{0}, x_{0}\right) $. Tenemos que 
\[y'\left(t\right)=\frac{\partial \varphi}{\partial t}\left(t; t_{0}, x_{0}\right)=f\left(t, y\left(t\right)\right) \quad \text{y} \quad y\left(s\right)=\varphi\left(s; t_{0}, x_{0}\right)=x^{*} .\]
Por otro lado, 
\[z'\left(t\right)=\frac{\partial \varphi}{\partial t}\left(t; s, \varphi\left(s; t_{0}, x_{0}\right)\right)=f\left(t, z\left(t\right)\right) .\]
Además, 
\[z\left(s\right)=\varphi\left(s; s, \varphi\left(s; t_{0}, x_{0}\right)\right)=\varphi\left(s; t_{0}, x_{0}\right)=x^{*} .\]
Tanto $\displaystyle y $ como $\displaystyle z $ son solución del mismo PVI y por unicidad de solución debe ser que $\displaystyle y\left(t\right)=z\left(t\right) $. 
\end{proof}
\section{Continuidad del flujo}
\begin{prop}[Dependencia respecto a $\displaystyle x_0 $]
Sea $\displaystyle f : D \subset \R \times \R^{n} \to \R^{n} $ continua y localmente Lipschitz respecto a la segunda variable y $\displaystyle \left(t_{0}, y_{0}\right), \left(t_{0}, z_{0}\right)\in D $ con
\[t_{0} \in I = \left[a,b\right] \subset\left(\beta_{-}\left(t_{0}, y_{0}\right), \beta_{+ }\left(t_{0}, y_{0}\right)\right)\cap \left(\beta_{-}\left(t_{0}, z_{0}\right), \beta_{+ }\left(t_{0}, z_{0}\right)\right) \]
y $\displaystyle K \subset D $ compacto tal que $\displaystyle \left(t, \varphi\left(t; t_{0}, y_{0}\right)\right), \left(t, \varphi\left(t; t_{0}, z_{0}\right)\right)\in K $, $\displaystyle \forall t \in I $. 
Entonces, 
\[ \|\varphi\left(t; t_{0}, y_{0}\right)-\varphi\left(t; t_{0}, z_{0}\right)\|\leq \|y_{0}-z_{0}\|e^{L \left|t-t_{0}\right|} , \; \forall t \in I\]
con $\displaystyle L \geq 0 $ constante de Lipschitz de $\displaystyle f $ en $\displaystyle K $. 
\end{prop}
\begin{proof}
Denominamos 
\[y\left(t\right):=\varphi\left(t; t_{0}, y_{0}\right)\quad \text{y} \quad z\left(t\right):=\varphi\left(t; t_{0}, z_{0}\right) .\]
Satisfacen que 
\[y\left(t\right)=y_{0} + \int^{t}_{t_{0}} f\left(s, y\left(s\right)\right) \; d s \quad \text{y} \quad z\left(t\right)=z_{0}+\int^{t}_{t_{0}} f\left(s, z\left(s\right)\right) \; d s .\]
Restando ambas igualdades tenemos que 
\[
\begin{split}
	\left\|y\left(t\right)-z\left(t\right)\right\| = & \left\|y_{0}-z_{0} + \int^{t}_{t_{0}} f\left(s,y\left(s\right)\right)-f\left(s,z\left(s\right)\right) \; d s\right\|\\
	\leq & \left\|y_{0}-z_{0}\right\| + \left\|\int^{t}_{t_{0}} f\left(s, y\left(s\right)\right)-f\left(s, z\left(s\right)\right) \; d s\right\|\\
	\leq & \|y_{0}-x_{0}\| + L \left\|\int^{t}_{t_{0}} \left\|y\left(s\right)-z\left(s\right)\right\| \; d s\right\|.
\end{split}
\]
Por la desigualdad de Grönwall tenemos que si tomamos $\displaystyle u\left(t\right)= \left\|y\left(t\right)-z\left(t\right)\right\| $, $\displaystyle A = \|y_{0}-z_{0}\| $ y $\displaystyle B = L $, tenemos que 
\[\|y\left(t\right)-z\left(t\right)\| \leq \|y_{0}-x_{0}\|e^{L \left|t-t_{0}\right|} .\]
\end{proof}
\begin{prop}[Dependencia respecto a $\displaystyle t_0 $]
Sea $\displaystyle f : D \subset \R \times \R^{n} \to \R^{n} $ continua y localmente Lipschitz respecto de la segunda variable, $\displaystyle \left(t_{0}, x_{0}\right), \left(t_{1}, x_{0}\right)\in D $, 
\[t_{0}, t_{1} \in I = \left[a,b\right] \subset \left(\beta_{-}\left(t_{0}, x_{0}\right), \beta_{+ }\left(t_{0}, x_{0}\right)\right) \cap \left(\beta_{- }\left(t_{1}, x_{0}\right), \beta_{+ }\left(t_{1},x_{0}\right)\right) \]
y $\displaystyle K \subset D $ compacto tal que $\displaystyle \left(t, \varphi\left(t; t_{0}, x_{0}\right)\right), \left(t, \varphi\left(t; t_{1}, x_{0}\right)\right)\in K $, $\displaystyle \forall t \in I $. Entonces,
\[\|\varphi\left(t; t_{0}, x_{0}\right)-\varphi\left(t; t_{1}, x_{0}\right)\| \leq M \left|t_{0}-t_{1}\right|e^{L \left|t-t_{1}\right|}, \; \forall t \in I ,\]
con $\displaystyle L \geq 0 $ constante de Lipschitz de $\displaystyle f $ en $\displaystyle K $ y $\displaystyle M $ máximo de $\displaystyle \left|f\right| $ en $\displaystyle K $.
\end{prop}
\begin{proof}
Definimos
\[y\left(t\right)=\varphi\left(t; t_{0}, x_{0}\right)\quad \text{y} \quad z\left(t\right)=\varphi\left(t; t_{1}, x_{0}\right) .\]
Tenemos que satisfacen 
\[ y\left(t\right)=x_{0} + \int^{t}_{t_{0}} f\left(s,y\left(s\right)\right) \; d s \quad \text{y} \quad z\left(t\right)=x_{0} +\int^{t}_{t_{1}} f\left(s, z\left(s\right)\right) \; d s .\]
Restando ambas expresiones obtenemos, 
\[
\begin{split}
	\left\|y\left(t\right)-z\left(t\right)\right\|= & \left\|\int^{t}_{t_{0}} f\left(s, y\left(s\right)\right) \; d s - \int^{t}_{t_{1}} f\left(s, z\left(s\right)\right) \; d s\right\| \\
	= & \left\|\int^{t_{1}}_{t_{0}} f\left(s, y\left(s\right)\right) \; d s + \int^{t}_{t_{1}} f\left(s, y\left(s\right)\right)-f\left(s,z\left(s\right)\right) \; d s\right\| \\
	\leq & \left\|\int^{t_{1}}_{t_{0}} f\left(s,y\left(s\right)\right) \; d s\right\| + \left\|\int^{t}_{t_{1}} f\left(s, y\left(s\right)\right)-f\left(s, z\left(s\right)\right) \; d s \right\| \\
	\leq & M \left|t_{1}-t_{0}\right| + L \left\|\int^{t}_{t_{1}} \|y\left(s\right)-z\left(s\right)\| \; d s\right\|.
\end{split}
\]
Aplicando la desigualdad de Grönwall, 
\[\|y\left(t\right)-z\left(t\right)\| \leq M \left|t_{1}-t_{0}\right|e^{L \left|t-t_{1}\right|} .\]
\end{proof}
\begin{theorem}[Continuidad del flujo]
	Sea $\displaystyle f : D \subset \R \times \R^{n} \to \R^{n} $ continua y localmente Lipschitz respecto de la segunda variable. Entonces $\displaystyle \varphi : \tilde{D} \subset \R \times \R \times \R^{n} \to \R^{n} $ es una función continua. 
\end{theorem}
\begin{proof}
	Fijamos $\displaystyle \left(\tau_{0}, t_{0}, x_{0}\right)\in \tilde{D} $ y veamos que $\displaystyle \varphi $ es continua en ese punto. Sea $\displaystyle \left[a,b\right] \subset I_{\left(t_{0}, x_{0}\right)} $ tal que $\displaystyle \tau_{0}, t_{0} \in \left[a,b\right]  $. Tomamos $\displaystyle \epsilon > 0 $ pequeño tal que el compacto
	\[K_{\epsilon }:= \left\{ \left(t,x\right) \; : \; \|x-\varphi\left(t; t_{0}, x_{0}\right)\| \leq 2\epsilon , t \in \left[a,b\right] \right\}  \subset D.\]
Elegimos 
\[\delta_{1} = \frac{\epsilon }{1+M}\frac{1}{1 + e^{L\left(b-a\right)}} > 0 ,\]
con $\displaystyle M = \max_{K_{\epsilon }} \left|f\right| $ y $\displaystyle L $ es la constante de Lipschitz de $\displaystyle f $ en $\displaystyle K_{\epsilon } $. Distinguimos varios pasos en lo que queda de demostración.
\begin{itemize}
	\item Si $\displaystyle \left|t_{1}-t_{0}\right|<\delta _{1} $ y $\displaystyle \|x_{1}-x_{0}\| < \delta _{1} $, queremos ver que $\displaystyle \left(t, \varphi\left(t; t_{1}, x_{1}\right)\right) \in K_{\epsilon } $, $\displaystyle \forall t \in \left[a,b\right]  $. Definimos $\displaystyle \left[c,d\right] \subset \left[a,b\right]  $ tal que sea el intervalo más grande que contiene a $\displaystyle t_{1} $ tal que $\displaystyle \left(t, \varphi\left(t; t_{1}, x_{1}\right)\right)\in K_{\epsilon } $. Podemos comprobar que $\displaystyle t_{1}\in \left[c,d\right]  $. 
Llamamos $\displaystyle z\left(t\right)=\varphi\left(t; t_{1}, x_{1}\right) $. Entonces tenemos que 
\[z\left(t\right)=x_{1} + \int^{t}_{t_{1}} f\left(s, z\left(s\right)\right) \; d s, \; \forall t \in I_{\left(t_{1}, x_{1}\right)} .\]
En particular, es cierto en $\displaystyle \left[c,d\right]  $. De esta forma para cada $\displaystyle t \in \left[c,d\right]  $ se tiene que 
\[
\begin{split}
	\|z\left(t\right)-\varphi\left(t; t_{0}, x_{0}\right)\| = & \left\|x_{1}-\varphi\left(t; t_{0}, x_{0}\right)+\int^{t}_{t_{1}} f\left(s,z\left(s\right)\right) \; d s\right\| \\
	\leq & \|x_{1}-x_{0}\| + \left\|\int^{t}_{t_{1}} f\left(s,z\left(s\right)\right) \; d s-\int^{t}_{t_{0}} f\left(s, \varphi\left(t; t_{0}, x_{0}\right)\right) \; d s\right\| \\
	\leq & \|x_{1}-x_{0}\| + \left\|\int^{t_{0}}_{t_{1}} \|f\left(s, z\left(s\right)\right)\| \; d s\right\| + \left\|\int^{t}_{t_{0}} \|f\left(s, z\left(s\right)\right)-f\left(s, \varphi\left(s; t_{0}, x_{0}\right)\right)\| \; d s\right\|\\
	\leq & \|x_{1}-x_{0}\| + M \left|t_{1}-t_{0}\right| + L\left\|\int^{t_{1}}_{t_{0}} \|z\left(s\right)-\varphi\left(s; t_{0}, x_{0}\right) \; d s\right\|
\end{split}
\]
Aplicando la desigualdad de Grönwall nos queda que para todo $\displaystyle t \in \left[c,d\right]  $,
\[\|z\left(t\right)-\varphi\left(t; t_{0}, x_{0}\right)\| \leq \left(\|x_{1} -x_{0}\| + M \left|t_{0}-t_{1}\right|\right)e^{L \left|t-t_{0}\right|} \leq \delta_{1}\left(1+M\right)e^{L\left(b-a\right)}<\epsilon   .\]
Si $\displaystyle \left[c,d\right] \neq \left[a,b\right]  $ tenemos que por continuidad, en $\displaystyle t = c $ o $\displaystyle t = d $, 
\[ \|z\left(t\right)-\varphi\left(t; t_{0}, x_{0}\right)\|\geq 2\epsilon  ,\]
que es una contradicción con la desigualdad que hemos obtenido anteriormente. Así, debe ser que $\displaystyle \left[c,d\right] =\left[a,b\right]  $. 
\item Tomamos $\displaystyle \delta \in \left(0, \delta _{1}\right) $ de forma que $\displaystyle \left(t_{0}-\delta, t_{0} + \delta \right) $ y $\displaystyle \left(\tau_{0}-\delta, \tau_{0} + \delta \right) $ estén contenidos en $\displaystyle \left[a,b\right]  $.
	Si $\displaystyle \|\left(\tau_{0}, t_{0}, x_{0}\right)-\left(\tau_{1}, t_{1}, x_{1}\right)\| < \delta  $, se tiene que 
	\[
	\begin{split}
		\|\varphi\left(\tau_{0}, t_{0}, x_{0}\right)-\varphi\left(\tau_{1}, t_{1}, x_{1}\right)\| \leq & \|\varphi\left(\tau_{0}, t_{0}, x_{0}\right)-\varphi\left(\tau_{1}, t_{0}, x_{0}\right)\| + \|\varphi\left(\tau_{1}, t_{0}, x_{0}\right)-\varphi\left(\tau_{1}, t_{1}, x_{1}\right)\| \\
		= & \left\|\int^{\tau_{1}}_{\tau_{0}} f\left(s, \varphi\left(s, t_{0}, x_{0}\right)\right) \; d s\right\| + M \left|t_{0}-t_{1}\right|e^{L \left(b-a\right)} + \|x_{0}-x_{1}\|e^{L\left(b-a\right)} \\
		\leq & M \|\tau_{1} -\tau_{0}\| + M \left|t_{0}-t_{1}\right|e^{L\left(b-a\right)} + \|x_{0}-x_{1}\|e^{L\left(b-a\right)} < \epsilon .
	\end{split}
	\]
\end{itemize}
\end{proof}

\begin{observation}
Hemos demostrado que el flujo no es sólo continuo sino que también es localmente Lipschitz. 
\end{observation}
\section{Dependencia respecto a parámetros}
\begin{prop}
Todo PVI en $\displaystyle \R^{n} $ con $\displaystyle m $ parámetros es equivalente a un PVI sin parámetros en $\displaystyle \R^{n + m} $. 
\end{prop}
\begin{proof}
Tenemos el PVI
\[
\begin{cases}
x' = f\left(t,x,\lambda \right) \\
x\left(t_{0}\right)=x_{0}
\end{cases}
, \]
con $\displaystyle \lambda \in \R^{m} $. Definimos
\[y\left(t\right)=\begin{pmatrix} x\left(t\right)\\ \lambda  \end{pmatrix}\in \R^{n+m}, \; \forall t .\]
Tendremos que 
\[y'\left(t\right)=\begin{pmatrix} x'\left(t\right) \\ 0 \end{pmatrix}=\begin{pmatrix} f\left(t, x\left(t\right), \lambda \right) \\ 0 \end{pmatrix}=\begin{pmatrix} f\left(t, y\left(t\right)\right) \\ 0 \end{pmatrix} .\]
Además, $\displaystyle y\left(t_{0}\right) = \left(x_{0}, \lambda \right) $. De esta forma obtenemos el PVI
\[
\begin{cases}
y'\left(t\right)=F\left(t, y\left(t\right)\right) \\
y\left(t_{0}\right)=\left(x_{0}, \lambda \right)
\end{cases}
,\]
con $\displaystyle F\left(t, y\left(t\right)\right):=\left(f\left(t, y\left(t\right)\right), 0\right) $. Es un PVI en $\displaystyle \R^{n+ m} $ puesto que $\displaystyle F : \R \times \left(\R^{n} \times \R^{m}\right)\to \R^{n+m} $. 
\end{proof}
\section{Dependencia diferenciable}
\begin{theorem}
	Si $\displaystyle f \in \mathcal{C}^{k}\left(D, \R^{n}\right) $ con $\displaystyle k \in \N $, entonces $\displaystyle \varphi \in \mathcal{C}^{k}\left(\tilde{D}, \R^{n}\right) $. 
\end{theorem}
\begin{proof}
La demostración es bastante laboriosa por lo que sólo daremos una idea de la demostración. Basta con demostrarlo para $\displaystyle k = 1 $ y generalizar usando inducción. Tenemos que ver que existen $\displaystyle \frac{\partial \varphi}{\partial t} $, $\displaystyle \frac{\partial \varphi}{\partial t_{0}} $ y $\displaystyle D\varphi\left(x_{0}\right) $ y son continuas. 
\begin{itemize}
\item En primer lugar, 
\[
\frac{\partial \varphi}{\partial t}=f\left(t, \varphi\left(t; t_{0}, x_{0}\right)\right) .
\]
Como $\displaystyle f $ es continua y localmente Lipschitz, en particular tendremos que $\displaystyle \varphi $ es continua por el teorema de continuidad. Por tanto, $\displaystyle \frac{\partial \varphi}{\partial t} $ es continua por ser composición de funciones continuas.
\item Supongamos que $\displaystyle D_{x_{0}}\varphi $ existe \footnote{No vamos a demostrar rigurosamente que existe porque requiere el uso de aproximaciones sucesivas.}, entonces
\[\frac{\partial }{\partial t}\left(D_{x_{0}}\varphi\right)=D_{x_{0}}\left(\frac{\partial \varphi}{\partial t}\right)=D_{x_{0}}\left(f\left(t; \varphi\left(t; t_{0}, x_{0}\right)\right)\right) = D_{x}f\left(f, \varphi\left(t; t_{0}, x_{0}\right)\right)D_{x_{0}}\varphi\left(t; t_{0}, x_{0}\right) .\]
Por la condición inicial $\displaystyle \varphi\left(t_{0}; t_{0}, x_{0}\right)=x_{0} $, tenemos que $\displaystyle D_{x_{0}}\varphi\left(t_{0}; t_{0}, x_{0}\right)=I $. La función matricial $\displaystyle D_{x_{0}}\varphi $ cumple que 
\[
\begin{cases}
\frac{\partial }{\partial t}\left(D_{x_{0}}\varphi\right)=D_{x}f\left(t, \varphi\right)D_{x_{0}}\varphi \\
D_{x_{0}}\varphi\left(t_{0}\right)= I
\end{cases}
.\]
Por tanto es solución de un PVI lineal. Como $\displaystyle D_{x}f\left(t, \varphi\right) $ es continua por ser composición de funciones continuas, del teorema de continuidad aplicado a sistemas lineales que la solución $\displaystyle D_{x_{0}}\varphi $ es continua.
\item Sabemos que $\displaystyle \varphi\left(t; t_{0} + h, x_{0}\right)=\varphi\left(t; t_{0}, \varphi\left(t_{0}; t_{0} + h, x_{0}\right)\right) $. Derivando respecto a $\displaystyle h $,
\[\frac{\partial\varphi}{\partial t_{0}}\left(t; t_{0} + h, x_{0}\right) = D_{x_{0}}\varphi\left(t; t_{0}, \varphi\left(t_{0}; t_{0} + h, x_{0}\right)\right)\frac{\partial \varphi}{\partial t_{0}}\left(t_{0}; t_{0} + h, x_{0}\right) .\]
Evaluando en $\displaystyle h = 0 $, 
\[\frac{\partial \varphi}{\partial t_{0}}\left(t; t_{0}, x_{0}\right)=D_{x_{0}}\varphi\left(t; t_{0}, x_{0}\right)\frac{\partial \varphi}{\partial t_{0}}\left(t_{0}; t_{0}, x_{0}\right) .\]
Por otro lado, $\displaystyle \varphi\left(t_{0}; t_{0}, x_{0}\right)=x_{0} $ y derivando respecto a $\displaystyle t_{0} $,
\[0 = \frac{\partial \varphi}{\partial t}\left(t_{0}; t_{0}, x_{0}\right) + \frac{\partial \varphi}{\partial t_{0}}\left(t_{0}; t_{0}, x_{0}\right)=f\left(t_{0}, \varphi\left(t_{0}; t_{0}, x_{0}\right)\right) + \frac{\partial \varphi}{\partial t_{0}}\left(t_{0}; t_{0}, x_{0}\right)=f\left(t_{0}, x_{0}\right) + \frac{\partial \varphi}{\partial t_{0}}\left(t_{0}; t_{0}, x_{0}\right).\]
Sustituyendo en la ecuación anterior,
\[\frac{\partial \varphi}{\partial t_{0}}\left(t; t_{0}, x_{0}\right)=-D_{x_{0}}\varphi\left(t; t_{0}, x_{0}\right)f\left(t_{0}, x_{0}\right) ,\]
que es una función continua por ser producto de funciones continuas. 
\end{itemize}
\end{proof}

